\section*{Capítulo \ref{ch:b1:nucleophilic-substitution} --- Substituição nucleofílica}

\begin{solution}{exo:b1:nucleophilic-substitution:1}
\ce{CH3CH2OH + Br-}; \ce{CH3OCH2CH3 + I-}; \ce{CH3CH2CH2CN + Cl-};
\ce{CH3NH3+ + Br-}. \omterm{def:b1:nucleophilic-substitution:substitution}{Substratos}: os haloalcanos; \omterm{def:b1:electronic-effects:nucleophile}{nucleófilos}:
\ce{OH-}, \ce{CH3CH2O-}, \ce{CN-}, \ce{NH3}; \omterm{def:b1:electronic-effects:nucleophile}{grupos de saída}: os íons haleto.
\end{solution}

\begin{solution}{exo:b1:nucleophilic-substitution:2}
Segunda ordem: triplicada; ambas reduzidas à metade, dividida por 4.
Primeira ordem só em RY: inalterada; reduzida à metade.
\end{solution}

\begin{solution}{exo:b1:nucleophilic-substitution:3}
SN2 (primário, bom \omterm{def:b1:electronic-effects:nucleophile}{nucleófilo}, aprótico); SN1 (terciário, água); SN2
(metila, \omterm{def:b1:electronic-effects:nucleophile}{nucleófilo} forte); \omterm{def:b1:nucleophilic-substitution:sn1}{solvólise} SN1 (secundário, \omterm{def:b1:electronic-effects:nucleophile}{nucleófilo} fraco
e neutro, \omterm{def:b1:intermolecular-forces-solvents:solvent-classes}{solvente prótico}), possivelmente com um pouco de SN2.
\end{solution}

\begin{solution}{exo:b1:nucleophilic-substitution:4}
O cianeto toma o lugar oposto ao bromo e os três outros grupos se viram.
No produto, o grupo CN é o de maior prioridade, como era o Br, e as
outras prioridades não mudam: o arranjo invertido tem o descritor
oposto, \cip{S}-2-metilbutanonitrila.
\end{solution}

\begin{solution}{exo:b1:nucleophilic-substitution:5}
1. \omterm{def:b1:electronic-effects:cleavage}{Heterólise}: seta da ligação C--Br para o Br, dando \ce{(CH3)3C+} e
\ce{Br-} (lenta). 2. Um \omterm{def:b1:lewis-resonance-vsepr:lewis}{par não ligante} da água para o carbono
catiônico: \ce{(CH3)3C-OH2+}. 3. Um \omterm{def:b1:lewis-resonance-vsepr:lewis}{par não ligante} de outra molécula de
água retira um próton do grupo \ce{OH2+}, com o par O--H voltando ao
oxigênio. Só a primeira etapa, lenta, decide a velocidade; a água, o
solvente, está em tal excesso que sua concentração não muda.
\end{solution}

\begin{solution}{exo:b1:nucleophilic-substitution:6}
O carbono que leva o Br é primário, mas ao lado dele um grupo terc-butila
ocupa o espaço atrás da ligação C--Br: o \omterm{def:b1:elementary-steps:energy-profile}{estado de transição} da SN2 fica
muito congestionado. A SN1 exigiria um \omterm{def:b1:electronic-effects:intermediates}{carbocátion} primário, instável
demais.
\end{solution}

\begin{solution}{exo:b1:nucleophilic-substitution:7}
O iodeto substitui o iodeto por SN2, a cada vez com inversão: \cip{S}
vira \cip{R} até que os dois estejam igualmente presentes, uma \omterm{def:b1:stereochemistry-in-depth:enantiomers}{mistura
racêmica}. Cada substituição transforma uma molécula \cip{S} em uma
\cip{R}: a rotação cai o equivalente a duas moléculas, de modo que a
rotação diminui duas vezes mais depressa do que o \omterm{def:b1:nucleophilic-substitution:substitution}{substrato} é
substituído.
\end{solution}

\begin{solution}{exo:b1:nucleophilic-substitution:8}
SN2: bromometano $>$ 1-bromopropano $>$ 2-bromopropano $>$
2-bromo-2-metilpropano (congestionamento). SN1: a ordem inversa
(estabilidade do \omterm{def:b1:electronic-effects:intermediates}{carbocátion}), com o bromometano e o 1-bromopropano quase
sem reagir.
\end{solution}

\begin{solution}{exo:b1:nucleophilic-substitution:9}
$x_{\text{inv}} - x_{\text{ret}} = 0.12$ e $x_{\text{inv}} + x_{\text{ret}} = 1$:
\qty{56}{\percent} invertido, \qty{44}{\percent} retido. Sobretudo
\omterm{def:b1:nucleophilic-substitution:sn1}{racemização} SN1, com um ligeiro excesso de inversão: o brometo que sai
ainda protege o seu lado quando a água chega.
\end{solution}

\begin{solution}{exo:b1:nucleophilic-substitution:10}
$v = k_1k_2[\mathrm{RY}][\mathrm{Nu}]/(k_{-1}[\mathrm{Y^-}] + k_2[\mathrm{Nu}])$
(\cref{prop:b1:nucleophilic-substitution:sn1-law}). Adicionar \ce{Y-}
aumenta o denominador: mais \omterm{def:b1:electronic-effects:intermediates}{carbocátions} voltam ao \omterm{def:b1:nucleophilic-substitution:substitution}{substrato} em vez de
serem capturados, e a velocidade cai. Desprezível quando
$k_2[\mathrm{Nu}] \gg k_{-1}[\mathrm{Y^-}]$, como quando o \omterm{def:b1:electronic-effects:nucleophile}{nucleófilo} é o
solvente.
\end{solution}

\begin{solution}{exo:b1:nucleophilic-substitution:11}
Concentrações iguais: $-\mathrm{d}c/\mathrm{d}t = kc^2$, $1/c = 1/C_0 + kt$,
$t_{1/2} = 1/(kC_0) = \qty{5.0e5}{s}$ (cerca de seis dias). Com um
excesso de cinquenta vezes de hidróxido, $[\ce{OH-}] \approx 50C_0$ é
constante: $c = C_0
e^{-k't}$ com $k' = 50kC_0 = \qty{1.0e-4}{s^{-1}}$, $t_{1/2} = \ln 2/k' =
\qty{6.9e3}{s}$, cerca de duas horas.
\end{solution}

\begin{solution}{exo:b1:nucleophilic-substitution:12}
O brometo teria de expulsar \ce{OH-}, uma \omterm{def:b1:predominant-reaction:strong-acid}{base forte} e um mau \omterm{def:b1:electronic-effects:nucleophile}{grupo de
saída}: nenhuma reação. No HBr concentrado, o OH é protonado a
\ce{-OH2+}, cujo \omterm{def:b1:electronic-effects:nucleophile}{grupo de saída} é a água. Álcool terciário: a água sai,
dando o \omterm{def:b1:electronic-effects:intermediates}{carbocátion}, capturado por \ce{Br-} (SN1). Álcool primário:
\ce{Br-} ataca o carbono de \ce{R-OH2+} por trás enquanto a água sai
(SN2).
\end{solution}

\begin{solution}{pb:b1:nucleophilic-substitution:1}
\textbf{1.} Cada molécula que reage dá um \ce{H+} e um \ce{Cl-}; a
condutividade é a soma das contribuições iônicas, proporcional à
quantidade de íons.
\textbf{2.} $[\mathrm{RCl}]/[\mathrm{RCl}]_0 = (\kappa_\infty -
\kappa)/\kappa_\infty$.
\textbf{3.} $\ln[\mathrm{RCl}] = \ln[\mathrm{RCl}]_0 - kt$: trace
$\ln(\kappa_\infty - \kappa)$ em função de $t$.
\textbf{4.} 0.693, 0.603, 0.513, 0.423, 0.333, 0.153, $-0.117$, $-0.387$.
\textbf{5.} Sim, inclinação de $-0.0180$ por minuto o tempo todo:
primeira ordem.
\textbf{6.} Qualquer ordem em relação à água: sua concentração é
constante (degenerescência da ordem).
\textbf{7.} $k = 0.0180/60 = \qty{3.0e-4}{s^{-1}}$.
\textbf{8.} $t_{1/2} = \ln 2/k = \qty{2.3e3}{s}$, \qty{38.5}{min}.
\textbf{9.} $\ln 10/k = \qty{7.7e3}{s}$, \qty{128}{min}.
\textbf{10.} $\ln(\kappa_\infty - \kappa)$: 0.693 em $t = 0$, 0.109 em
\qty{10}{min}, e assim por diante: inclinação de $-0.0584$ por minuto,
$k = \qty{9.7e-4}{s^{-1}}$.
\textbf{11.} $9.7/3.0 = 3.2$.
\textbf{12.} Em \qty{30}{min}: $2.000(1 - e^{-0.0584 \times 30}) = 1.653$,
medido 1.654.
\textbf{13.} SN1: \omterm{def:b1:nucleophilic-substitution:substitution}{substrato} terciário, \omterm{def:b1:intermolecular-forces-solvents:solvent-classes}{solvente prótico}, \omterm{def:b1:electronic-effects:nucleophile}{nucleófilo}
neutro.
\textbf{14.} Como no \cref{exo:b1:nucleophilic-substitution:5}, com Cl.
\textbf{15.} A ionização lenta só envolve o \omterm{def:b1:nucleophilic-substitution:substitution}{substrato}: $v =
k[\mathrm{RCl}]$.
\textbf{16.} O hidróxido só poderia agir na etapa rápida de captura,
depois da ionização determinante; a velocidade não depende dele.
\textbf{17.} O \omterm{def:b1:electronic-effects:intermediates}{carbocátion} é plano e \omterm{def:b1:stereochemistry-in-depth:stereocentre}{aquiral}; a água se adiciona
igualmente às duas faces: álcool racêmico.
\textbf{18.} Uma primeira barreira alta (ionização) até o \omterm{def:b1:electronic-effects:intermediates}{carbocátion},
em um poço raso, e depois uma barreira baixa para a captura: a primeira
etapa é a determinante.
\textbf{19.} Ele formaria um \omterm{def:b1:electronic-effects:intermediates}{carbocátion} primário (instável demais), e a
água é um \omterm{def:b1:electronic-effects:nucleophile}{nucleófilo} fraco demais para uma SN2 rápida.
\textbf{20.} $k = A\,e^{-E_a/RT}$.
\textbf{21.} $E_a = R\ln(k_2/k_1)/(1/T_1 - 1/T_2)$.
\textbf{22.} $E_a = 8.314 \times \ln(9.74/3.00)/(1/298.15 - 1/308.15) =
\qty{9.0e4}{J/mol}$.
\textbf{23.} A altura da barreira da ionização, a \omterm{def:b1:elementary-steps:rds}{etapa determinante}.
\textbf{24.} $k = 3.0 \times 10^{-4} \exp\bigl(\frac{90\,000}{8.314}
(\frac{1}{298.15} - \frac{1}{318.15})\bigr) = \qty{2.9e-3}{s^{-1}}$.
\textbf{25.} \textbf{$E_a = \qty{90}{kJ/mol}$}.
\end{solution}
